% Options for packages loaded elsewhere
\PassOptionsToPackage{unicode}{hyperref}
\PassOptionsToPackage{hyphens}{url}
%
\documentclass[
]{article}
\usepackage{amsmath,amssymb}
\usepackage{iftex}
\ifPDFTeX
  \usepackage[T1]{fontenc}
  \usepackage[utf8]{inputenc}
  \usepackage{textcomp} % provide euro and other symbols
\else % if luatex or xetex
  \usepackage{unicode-math} % this also loads fontspec
  \defaultfontfeatures{Scale=MatchLowercase}
  \defaultfontfeatures[\rmfamily]{Ligatures=TeX,Scale=1}
\fi
\usepackage{lmodern}
\ifPDFTeX\else
  % xetex/luatex font selection
\fi
% Use upquote if available, for straight quotes in verbatim environments
\IfFileExists{upquote.sty}{\usepackage{upquote}}{}
\IfFileExists{microtype.sty}{% use microtype if available
  \usepackage[]{microtype}
  \UseMicrotypeSet[protrusion]{basicmath} % disable protrusion for tt fonts
}{}
\makeatletter
\@ifundefined{KOMAClassName}{% if non-KOMA class
  \IfFileExists{parskip.sty}{%
    \usepackage{parskip}
  }{% else
    \setlength{\parindent}{0pt}
    \setlength{\parskip}{6pt plus 2pt minus 1pt}}
}{% if KOMA class
  \KOMAoptions{parskip=half}}
\makeatother
\usepackage{xcolor}
\usepackage[margin=1in]{geometry}
\usepackage{color}
\usepackage{fancyvrb}
\newcommand{\VerbBar}{|}
\newcommand{\VERB}{\Verb[commandchars=\\\{\}]}
\DefineVerbatimEnvironment{Highlighting}{Verbatim}{commandchars=\\\{\}}
% Add ',fontsize=\small' for more characters per line
\usepackage{framed}
\definecolor{shadecolor}{RGB}{248,248,248}
\newenvironment{Shaded}{\begin{snugshade}}{\end{snugshade}}
\newcommand{\AlertTok}[1]{\textcolor[rgb]{0.94,0.16,0.16}{#1}}
\newcommand{\AnnotationTok}[1]{\textcolor[rgb]{0.56,0.35,0.01}{\textbf{\textit{#1}}}}
\newcommand{\AttributeTok}[1]{\textcolor[rgb]{0.13,0.29,0.53}{#1}}
\newcommand{\BaseNTok}[1]{\textcolor[rgb]{0.00,0.00,0.81}{#1}}
\newcommand{\BuiltInTok}[1]{#1}
\newcommand{\CharTok}[1]{\textcolor[rgb]{0.31,0.60,0.02}{#1}}
\newcommand{\CommentTok}[1]{\textcolor[rgb]{0.56,0.35,0.01}{\textit{#1}}}
\newcommand{\CommentVarTok}[1]{\textcolor[rgb]{0.56,0.35,0.01}{\textbf{\textit{#1}}}}
\newcommand{\ConstantTok}[1]{\textcolor[rgb]{0.56,0.35,0.01}{#1}}
\newcommand{\ControlFlowTok}[1]{\textcolor[rgb]{0.13,0.29,0.53}{\textbf{#1}}}
\newcommand{\DataTypeTok}[1]{\textcolor[rgb]{0.13,0.29,0.53}{#1}}
\newcommand{\DecValTok}[1]{\textcolor[rgb]{0.00,0.00,0.81}{#1}}
\newcommand{\DocumentationTok}[1]{\textcolor[rgb]{0.56,0.35,0.01}{\textbf{\textit{#1}}}}
\newcommand{\ErrorTok}[1]{\textcolor[rgb]{0.64,0.00,0.00}{\textbf{#1}}}
\newcommand{\ExtensionTok}[1]{#1}
\newcommand{\FloatTok}[1]{\textcolor[rgb]{0.00,0.00,0.81}{#1}}
\newcommand{\FunctionTok}[1]{\textcolor[rgb]{0.13,0.29,0.53}{\textbf{#1}}}
\newcommand{\ImportTok}[1]{#1}
\newcommand{\InformationTok}[1]{\textcolor[rgb]{0.56,0.35,0.01}{\textbf{\textit{#1}}}}
\newcommand{\KeywordTok}[1]{\textcolor[rgb]{0.13,0.29,0.53}{\textbf{#1}}}
\newcommand{\NormalTok}[1]{#1}
\newcommand{\OperatorTok}[1]{\textcolor[rgb]{0.81,0.36,0.00}{\textbf{#1}}}
\newcommand{\OtherTok}[1]{\textcolor[rgb]{0.56,0.35,0.01}{#1}}
\newcommand{\PreprocessorTok}[1]{\textcolor[rgb]{0.56,0.35,0.01}{\textit{#1}}}
\newcommand{\RegionMarkerTok}[1]{#1}
\newcommand{\SpecialCharTok}[1]{\textcolor[rgb]{0.81,0.36,0.00}{\textbf{#1}}}
\newcommand{\SpecialStringTok}[1]{\textcolor[rgb]{0.31,0.60,0.02}{#1}}
\newcommand{\StringTok}[1]{\textcolor[rgb]{0.31,0.60,0.02}{#1}}
\newcommand{\VariableTok}[1]{\textcolor[rgb]{0.00,0.00,0.00}{#1}}
\newcommand{\VerbatimStringTok}[1]{\textcolor[rgb]{0.31,0.60,0.02}{#1}}
\newcommand{\WarningTok}[1]{\textcolor[rgb]{0.56,0.35,0.01}{\textbf{\textit{#1}}}}
\usepackage{graphicx}
\makeatletter
\def\maxwidth{\ifdim\Gin@nat@width>\linewidth\linewidth\else\Gin@nat@width\fi}
\def\maxheight{\ifdim\Gin@nat@height>\textheight\textheight\else\Gin@nat@height\fi}
\makeatother
% Scale images if necessary, so that they will not overflow the page
% margins by default, and it is still possible to overwrite the defaults
% using explicit options in \includegraphics[width, height, ...]{}
\setkeys{Gin}{width=\maxwidth,height=\maxheight,keepaspectratio}
% Set default figure placement to htbp
\makeatletter
\def\fps@figure{htbp}
\makeatother
\setlength{\emergencystretch}{3em} % prevent overfull lines
\providecommand{\tightlist}{%
  \setlength{\itemsep}{0pt}\setlength{\parskip}{0pt}}
\setcounter{secnumdepth}{-\maxdimen} % remove section numbering
\ifLuaTeX
  \usepackage{selnolig}  % disable illegal ligatures
\fi
\usepackage{bookmark}
\IfFileExists{xurl.sty}{\usepackage{xurl}}{} % add URL line breaks if available
\urlstyle{same}
\hypersetup{
  pdftitle={HW 3},
  pdfauthor={Bryan Mui - UID 506021334},
  hidelinks,
  pdfcreator={LaTeX via pandoc}}

\title{HW 3}
\author{Bryan Mui - UID 506021334}
\date{2025-02-12}

\begin{document}
\maketitle

\subsection{\texorpdfstring{Things to ask
\n}{Things to ask }}\label{things-to-ask}

\begin{itemize}
\tightlist
\item
  1a, what is the correct kind of output
\item
  1b, how to do
\end{itemize}

\subsection{Section 1: Cook's relational
data}\label{section-1-cooks-relational-data}

library and read csv files:

\begin{verbatim}
## -- Attaching core tidyverse packages ------------------------ tidyverse 2.0.0 --
## v dplyr     1.1.4     v readr     2.1.5
## v forcats   1.0.0     v stringr   1.5.1
## v ggplot2   3.5.1     v tibble    3.2.1
## v lubridate 1.9.4     v tidyr     1.3.1
## v purrr     1.0.2     
## -- Conflicts ------------------------------------------ tidyverse_conflicts() --
## x dplyr::filter() masks stats::filter()
## x dplyr::lag()    masks stats::lag()
## i Use the conflicted package (<http://conflicted.r-lib.org/>) to force all conflicts to become errors
## Rows: 5 Columns: 2
## -- Column specification --------------------------------------------------------
## Delimiter: ","
## chr (2): name, Inventor
## 
## i Use `spec()` to retrieve the full column specification for this data.
## i Specify the column types or set `show_col_types = FALSE` to quiet this message.
## Rows: 22 Columns: 3
## -- Column specification --------------------------------------------------------
## Delimiter: ","
## chr (2): recipe, food_item
## dbl (1): weight (oz)
## 
## i Use `spec()` to retrieve the full column specification for this data.
## i Specify the column types or set `show_col_types = FALSE` to quiet this message.
## Rows: 25 Columns: 3
## -- Column specification --------------------------------------------------------
## Delimiter: ","
## chr (2): food_item, shop
## dbl (1): price (US dollars per lb)
## 
## i Use `spec()` to retrieve the full column specification for this data.
## i Specify the column types or set `show_col_types = FALSE` to quiet this message.
## Rows: 11 Columns: 3
## -- Column specification --------------------------------------------------------
## Delimiter: ","
## chr (2): item, type
## dbl (1): calories
## 
## i Use `spec()` to retrieve the full column specification for this data.
## i Specify the column types or set `show_col_types = FALSE` to quiet this message.
## Rows: 5 Columns: 2
## -- Column specification --------------------------------------------------------
## Delimiter: ","
## chr (2): name, Inventor
## 
## i Use `spec()` to retrieve the full column specification for this data.
## i Specify the column types or set `show_col_types = FALSE` to quiet this message.
## Rows: 22 Columns: 3
## -- Column specification --------------------------------------------------------
## Delimiter: ","
## chr (2): recipe, food_item
## dbl (1): weight (oz)
## 
## i Use `spec()` to retrieve the full column specification for this data.
## i Specify the column types or set `show_col_types = FALSE` to quiet this message.
## Rows: 25 Columns: 3
## -- Column specification --------------------------------------------------------
## Delimiter: ","
## chr (2): food_item, shop
## dbl (1): price (US dollars per lb)
## 
## i Use `spec()` to retrieve the full column specification for this data.
## i Specify the column types or set `show_col_types = FALSE` to quiet this message.
## Rows: 11 Columns: 3
## -- Column specification --------------------------------------------------------
## Delimiter: ","
## chr (2): item, type
## dbl (1): calories
## 
## i Use `spec()` to retrieve the full column specification for this data.
## i Specify the column types or set `show_col_types = FALSE` to quiet this message.
\end{verbatim}

\subsubsection{1a) Write R code to return the types of food and the
total calories for each type to make a turkey burger from the
recipe.}\label{a-write-r-code-to-return-the-types-of-food-and-the-total-calories-for-each-type-to-make-a-turkey-burger-from-the-recipe.}

\begin{Shaded}
\begin{Highlighting}[]
\NormalTok{a1}
\end{Highlighting}
\end{Shaded}

\begin{verbatim}
## # A tibble: 4 x 2
##   type          `sum(calories)`
##   <chr>                   <dbl>
## 1 Meat                     15  
## 2 Sauce                    35  
## 3 Vegetables                3.5
## 4 Wheat product            25
\end{verbatim}

\subsubsection{1b) Write R code to return all recipes containing both
bread and tomato products with their weights on these two
ingredients.}\label{b-write-r-code-to-return-all-recipes-containing-both-bread-and-tomato-products-with-their-weights-on-these-two-ingredients.}

\begin{Shaded}
\begin{Highlighting}[]
\NormalTok{b1}
\end{Highlighting}
\end{Shaded}

\begin{verbatim}
## # A tibble: 3 x 2
##   recipe        sum_bread_and_tomato
##   <chr>                        <dbl>
## 1 Beef Burger                     16
## 2 Sandwich                        17
## 3 Turkey Burger                   14
\end{verbatim}

\subsubsection{1c) Write R code to return the total calories of the beef
and turkey
burgers.}\label{c-write-r-code-to-return-the-total-calories-of-the-beef-and-turkey-burgers.}

\begin{Shaded}
\begin{Highlighting}[]
\NormalTok{c1}
\end{Highlighting}
\end{Shaded}

\begin{verbatim}
## # A tibble: 2 x 2
##   recipe        total_calories
##   <chr>                  <dbl>
## 1 Beef Burger             1317
## 2 Turkey Burger            777
\end{verbatim}

\subsubsection{1d) Write R code to return all vegetable items sold at
either `W-Mart' or `Food Warehouse' and display which shop offers the
lowest price for each
item.}\label{d-write-r-code-to-return-all-vegetable-items-sold-at-either-w-mart-or-food-warehouse-and-display-which-shop-offers-the-lowest-price-for-each-item.}

\begin{Shaded}
\begin{Highlighting}[]
\NormalTok{d1}
\end{Highlighting}
\end{Shaded}

\begin{verbatim}
## # A tibble: 5 x 3
##   food_item minprice shop          
##   <chr>        <dbl> <chr>         
## 1 Cabbage        1.8 Food warehouse
## 2 Onions         1   W-Mart        
## 3 Onions         1   Coco Mart     
## 4 Spinach        2.5 Food warehouse
## 5 Tomato         3   W-Mart
\end{verbatim}

\subsubsection{1e) Write R code to return the recipes with the total
calories that contain ``Wheat
product''.}\label{e-write-r-code-to-return-the-recipes-with-the-total-calories-that-contain-wheat-product.}

\begin{Shaded}
\begin{Highlighting}[]
\NormalTok{e1}
\end{Highlighting}
\end{Shaded}

\begin{verbatim}
## # A tibble: 4 x 2
##   recipe                  total_kcal
##   <chr>                        <dbl>
## 1 Beef Burger                   1317
## 2 Sandwich                       562
## 3 Spaghetti and Meatballs       1579
## 4 Turkey Burger                  777
\end{verbatim}

\subsection{Section 2: Regular Expression
Exercises}\label{section-2-regular-expression-exercises}

\subsubsection{Part 1: Find names.txt and write regex patterns to answer
following
questions.}\label{part-1-find-names.txt-and-write-regex-patterns-to-answer-following-questions.}

\subsubsection{1a) Find all usernames that contain at least one numeric
character.}\label{a-find-all-usernames-that-contain-at-least-one-numeric-character.}

\begin{Shaded}
\begin{Highlighting}[]
\NormalTok{txt }\OtherTok{\textless{}{-}} \FunctionTok{readLines}\NormalTok{(}\StringTok{"./names.txt"}\NormalTok{)}
\NormalTok{txt}
\end{Highlighting}
\end{Shaded}

\begin{verbatim}
##  [1] "Abc One23"         "Horrible turtle"   "Messsages"        
##  [4] "keep_it_simple"    "hello world!"      "Zoran D Wang"     
##  [7] "myUsernameis210"   "abc defg"          "asml"             
## [10] "john"              "Edward Lazowska"   "123fionaFog"      
## [13] "Red chihuahua5"    "1"                 "CLEAN"            
## [16] "+plus+"            "omaha poshy"       "0maha p0shy"      
## [19] "OMAHA POSHY"       "Deacon @ Stephens" "Ajay Ann Bryan"
\end{verbatim}

\begin{Shaded}
\begin{Highlighting}[]
\CommentTok{\# remember to store in R file!!!!!!!!!!!!!!!!!!!}
\NormalTok{pat\_1\_a }\OtherTok{\textless{}{-}} \StringTok{"}\SpecialCharTok{\textbackslash{}\textbackslash{}}\StringTok{d"}
\FunctionTok{str\_view}\NormalTok{(txt, pat\_1\_a)}
\end{Highlighting}
\end{Shaded}

\begin{verbatim}
##  [1] | Abc One<2><3>
##  [7] | myUsernameis<2><1><0>
## [12] | <1><2><3>fionaFog
## [13] | Red chihuahua<5>
## [14] | <1>
## [18] | <0>maha p<0>shy
\end{verbatim}

\subsubsection{1b) Find all usernames that are exactly four characters
long and consist only of alphabetic
characters.}\label{b-find-all-usernames-that-are-exactly-four-characters-long-and-consist-only-of-alphabetic-characters.}

\begin{Shaded}
\begin{Highlighting}[]
\NormalTok{pat\_1\_b }\OtherTok{\textless{}{-}} \StringTok{"}\SpecialCharTok{\textbackslash{}\textbackslash{}}\StringTok{A[[:alpha:]][[:alpha:]][[:alpha:]][[:alpha:]]}\SpecialCharTok{\textbackslash{}\textbackslash{}}\StringTok{Z"}
\FunctionTok{str\_view}\NormalTok{(txt, pat\_1\_b)}
\end{Highlighting}
\end{Shaded}

\begin{verbatim}
##  [9] | <asml>
## [10] | <john>
\end{verbatim}

\subsubsection{1c) Find all usernames following the conventional name
format, where the given name precedes the family name, with optional
additional names in-between. Each name segment is separated by a single
white space and should begin with an uppercase letter followed by one or
more lowercase letters. Note that intermediate names should also start
with an uppercase letter followed by zero or more lowercase
letters.}\label{c-find-all-usernames-following-the-conventional-name-format-where-the-given-name-precedes-the-family-name-with-optional-additional-names-in-between.-each-name-segment-is-separated-by-a-single-white-space-and-should-begin-with-an-uppercase-letter-followed-by-one-or-more-lowercase-letters.-note-that-intermediate-names-should-also-start-with-an-uppercase-letter-followed-by-zero-or-more-lowercase-letters.}

\begin{Shaded}
\begin{Highlighting}[]
\NormalTok{pat\_1\_c }\OtherTok{\textless{}{-}} \StringTok{"}\SpecialCharTok{\textbackslash{}\textbackslash{}}\StringTok{b[a{-}zA{-}z][a{-}zA{-}z][a{-}zA{-}z][a{-}zA{-}z]"}
\FunctionTok{str\_view}\NormalTok{(txt, pat\_1\_c)}
\end{Highlighting}
\end{Shaded}

\begin{verbatim}
##  [2] | <Horr>ible <turt>le
##  [3] | <Mess>sages
##  [4] | <keep>_it_simple
##  [5] | <hell>o <worl>d!
##  [6] | <Zora>n D <Wang>
##  [7] | <myUs>ernameis210
##  [8] | abc <defg>
##  [9] | <asml>
## [10] | <john>
## [11] | <Edwa>rd <Lazo>wska
## [13] | Red <chih>uahua5
## [15] | <CLEA>N
## [16] | +<plus>+
## [17] | <omah>a <posh>y
## [19] | <OMAH>A <POSH>Y
## [20] | <Deac>on @ <Step>hens
## [21] | <Ajay> Ann <Brya>n
\end{verbatim}

\subsubsection{Part 2: Find cards.txt and write the regex patterns to
identify all the valid card numbers with the correct formatting
rules:}\label{part-2-find-cards.txt-and-write-the-regex-patterns-to-identify-all-the-valid-card-numbers-with-the-correct-formatting-rules}

\begin{itemize}
\tightlist
\item
  A Master card number begins with a 5 and it is exactly 16 digits long.
\item
  A Visa card number begins with a 4 and it is between 13 and 16 digits
  long
\end{itemize}

\begin{Shaded}
\begin{Highlighting}[]
\NormalTok{cards }\OtherTok{\textless{}{-}} \FunctionTok{readLines}\NormalTok{(}\StringTok{"./cards.txt"}\NormalTok{)}
\NormalTok{cards}
\end{Highlighting}
\end{Shaded}

\begin{verbatim}
##  [1] "5123456789101112"         "4789 0123 8910 1112"     
##  [3] "4444 9321 1230 3"         "5315 4011 1721 51"       
##  [5] "4987 9381 2457"           "4891 0870 8908 70987"    
##  [7] "5234 4567 8910 1112"      "58907890782309171"       
##  [9] "3008 9078 1891 7890"      "5192 9295 91828818"      
## [11] "4182 2884 1232 9582 2182"
\end{verbatim}

\subsubsection{2a) Write a regex pattern to match valid Master card
number and print all the valid numbers, grouped into sets of 4 separated
by a single
space.}\label{a-write-a-regex-pattern-to-match-valid-master-card-number-and-print-all-the-valid-numbers-grouped-into-sets-of-4-separated-by-a-single-space.}

\begin{Shaded}
\begin{Highlighting}[]
\CommentTok{\# Valid master card numbers }
\NormalTok{pat\_2\_a }\OtherTok{\textless{}{-}} \StringTok{"5.\{15\}"}
\FunctionTok{str\_view}\NormalTok{(cards, pat\_2\_a)}
\end{Highlighting}
\end{Shaded}

\begin{verbatim}
##  [1] | <5123456789101112>
##  [4] | <5315 4011 1721 5>1
##  [7] | <5234 4567 8910 1>112
##  [8] | <5890789078230917>1
## [10] | <5192 9295 918288>18
\end{verbatim}

\subsubsection{2b) Write a regex pattern to match valid Visa card number
and print all the valid numbers, grouped into sets of 4 separated by a
single
space.}\label{b-write-a-regex-pattern-to-match-valid-visa-card-number-and-print-all-the-valid-numbers-grouped-into-sets-of-4-separated-by-a-single-space.}

\subsubsection{Part 3: Find the passwords.txt and write regex patterns
to answer the following
questions.}\label{part-3-find-the-passwords.txt-and-write-regex-patterns-to-answer-the-following-questions.}

\begin{Shaded}
\begin{Highlighting}[]
\NormalTok{passwords }\OtherTok{\textless{}{-}} \FunctionTok{readLines}\NormalTok{(}\StringTok{"./passwords.txt"}\NormalTok{)}
\end{Highlighting}
\end{Shaded}

\begin{verbatim}
## Warning in readLines("./passwords.txt"): incomplete final line found on
## './passwords.txt'
\end{verbatim}

\begin{Shaded}
\begin{Highlighting}[]
\NormalTok{passwords}
\end{Highlighting}
\end{Shaded}

\begin{verbatim}
##  [1] "1234567"        "12345678"       "Strings78"      "appleO07"      
##  [5] "1brownie"       "asdfjkl"        "?helloWorld"    "90095"         
##  [9] "glhf1234"       "H1234!!ap"      "789afk"         "alllowercase"  
## [13] "ALLUPPERCASE"   "missingNumbers"
\end{verbatim}

\subsubsection{3a) Write a regex pattern to identify the passwords that
satisfies the requirements
below.}\label{a-write-a-regex-pattern-to-identify-the-passwords-that-satisfies-the-requirements-below.}

\begin{itemize}
\tightlist
\item
  Minimum 8 characters
\item
  Must contain at least one letter
\item
  Must contain at least one digit
\end{itemize}

\subsubsection{3b) Write a regex pattern to identify the passwords that
satisfies the requirements
below}\label{b-write-a-regex-pattern-to-identify-the-passwords-that-satisfies-the-requirements-below}

\begin{itemize}
\tightlist
\item
  Minimum 8 characters
\item
  Must contain at least one uppercase character
\item
  Must contain at least one lowercase character
\item
  Must contain at least one digit
\item
  Must contain at least one punctuation character
\end{itemize}

\subsubsection{Part 4: Please load wordlists.RData to R and write
regular expression patterns which will match all of the values in x and
none of the values in
y.}\label{part-4-please-load-wordlists.rdata-to-r-and-write-regular-expression-patterns-which-will-match-all-of-the-values-in-x-and-none-of-the-values-in-y.}

\subsubsection{a) Ranges}\label{a-ranges}

\subsubsection{b) Backrefs}\label{b-backrefs}

\end{document}
